%%%%%%%%%%%%%%%%%%%%%%%%%%%%%%%%%%%%%%%%%%%%%%%%%%%%%%%%%%%%%%%%%%%%%%%%
%
% Two-page extended abstract
% From Human Videos to Dexterous Robot Capability workshop (CoRL 2026)
% Formatting: NVIDIA tech report.
%
% Limit: 2 pages of main text; references, acknowledgments and
% appendices do not count.
%
% Build:  pdflatex main && bibtex main && pdflatex main && pdflatex main
%
%%%%%%%%%%%%%%%%%%%%%%%%%%%%%%%%%%%%%%%%%%%%%%%%%%%%%%%%%%%%%%%%%%%%%%%%

\pdfoutput=1

%%% == IMPORTANT ==
%%% For anonymized submission (no authors, logo, copyright or
%%% acknowledgments), use this:
%\documentclass[10pt, logo, copyright, anonymous]{nvidiaextabstract}

%%% For camera-ready, use this:
\documentclass[10pt]{nvidiaextabstract}

%%% Load required packages here (many are loaded by the class already:
%%% amsmath, amssymb, graphicx, booktabs, xcolor, hyperref, natbib,
%%% balance, enumitem, caption).
\usepackage{microtype}
\usepackage{multirow}
\usepackage{subcaption}
\usepackage{xspace}

%%%%%%%%%%%%%%%%%%%%%%%%%%%%%%%%%%%%%%%%%%%%%%%%%%%%%%%%%%%%%%%%%%%%%%%%

%%% Venue: printed in the first-page footer. Comment out to drop it.
\venue{From Human Videos to Dexterous Robot Capability (CoRL 2026)}

%%% Logo in the top-left corner of page 1 (shown with the `logo' class
%%% option). Options in [] go to \includegraphics; default is the NVIDIA
%%% logo at width=90pt. The header is about 40pt tall, so keep height below that.
%\logo[height=28pt]{assets/mylogo.png}

%%% Printed instead of the authors when the `anonymous' option is set.
\submissionID{maniphys.ai}

\title{Per-Episode Policy Selection with an Object-Orientation Reward for Video-to-Robot Grounding}
\shorttitle{maniphys.ai: Track 2 Technical Report}

%%% Authors: to be filled in (names, order, affiliations, emails).
\author[1]{Author list to be provided}
\affil[1]{maniphysAI lab}

\begin{abstract}
\textbf{Abstract:}
V2D Track 2 asks a Unitree G1 with Dex3 hands to reproduce 30 whole-body human-object interactions in simulation from clean trajectories, noisy trajectories or one video.
The official FlashCHORD reward tracks the object position but not its orientation, so a policy can keep the object in place while it rotates away from the reference.
We add one orientation reward term, keep the rest of the official recipe, and accept a new policy for an episode only under swap rules fixed before each experiment.
On 12 tracked Tier 2 and Tier 3 episodes the term raised 4096-world ADD-AUC on 11, by 0.08 to 0.37.
Together with a clean-room Tier 2 denoiser and a monocular Tier 3 reconstruction aligned at frame 0, the entry ranks first on the Tier 3 SP-SR board (3 November 2026).
\end{abstract}

%%%%%%%%%%%%%%%%%%%%%%%%%%%%%%%%%%%%%%%%%%%%%%%%%%%%%%%%%%%%%%%%%%%%%%%%

%%% Author-defined commands.
\definecolor{NvidiaGreen}{RGB}{118, 185, 0}
\newcommand{\method}{\textsc{Method}\xspace}
\newcommand{\rtho}{\textsc{r3o}\xspace}

%%%%%%%%%%%%%%%%%%%%%%%%%%%%%%%%%%%%%%%%%%%%%%%%%%%%%%%%%%%%%%%%%%%%%%%%

\begin{document}

\maketitle

%%%%%%%%%%%%%%%%%%%%%%%%%%%%%%%%%%%%%%%%%%%%%%%%%%%%%%%%%%%%%%%%%%%%%%%%

\section{Introduction}
\label{sec:intro}

Track 2 asks for a robot policy that reproduces each of 30 human-object interaction episodes, such as righting a fallen bar stool or handing a basket from one hand to the other.
The input comes in three tiers: clean trajectories (Tier 1), trajectories with synthetic noise (Tier 2) and a single exocentric video (Tier 3).
The tasks are whole-body and bimanual: one policy has to walk, make contact and carry the object without an assist force at the end of training.

The official recipe trains FlashSAC, a soft actor-critic variant~\cite{haarnoja2018sac}, with a reward whose only object term is position.
ADD-AUC averages the distance of object model points, so an object that is held at the right place but turned is still scored as an error.
We therefore add a dense orientation term and treat the challenge as 90 tracking problems (30 episodes, 3 tiers), because one episode moves a board by only 1/30 and single-seed results vary widely.
Our contributions are:
\begin{itemize}
  \item an object-orientation reward term (Eq.~\ref{eq:ori}) that raised 4096-world ADD-AUC on 11 of 12 tracked Tier 2 and Tier 3 episodes, by 0.08 to 0.37;
  \item a per-episode selection procedure with swap rules fixed in advance, and an upload gate that discounts own-reference gains by a measured board transfer ratio of 0.21;
  \item a clean-room Tier 2 denoiser and a monocular Tier 3 reconstruction chain with frame-0 alignment.
\end{itemize}

\section{Approach}
\label{sec:approach}

\paragraph{Base recipe.}
Every policy is trained with the official FlashSAC configuration for 350\,M environment steps with 2048 parallel worlds and the official nine-stage curriculum, in which the object assist force decays to zero.
The replay buffer holds 10\,M transitions on 48\,GB GPUs and 4\,M on a 24\,GB GPU.
\emph{Variant A} changes two values: the object position reward width ($\sigma$ 0.2\,m $\rightarrow$ 0.05\,m) and the object orientation termination (1.5 $\rightarrow$ 0.6\,rad).

\paragraph{Object-orientation reward (\rtho).}
The official reward has twelve terms, and its only object term is on position; orientation enters only through a binary termination.
We add a thirteenth term
\begin{equation}
\label{eq:ori}
r_{\mathrm{ori}} = w \exp\!\left(-\theta_t^2 / \sigma^2\right),
\end{equation}
where $\theta_t$ is the geodesic angle between the simulated and the reference object orientation, with $w=1.0$ and $\sigma=0.4$\,rad in all curriculum stages.
The term gives a gradient on orientation before the termination fires.
It is set through two environment variables, so checkpoints keep the official configuration and the official evaluation code reads them unchanged.

\paragraph{Selection.}
Each candidate is evaluated in 4096 deterministic worlds against the reference of its own tier; Tier 2 and Tier 3 never use Tier 1 trajectories.
A candidate replaces the listed policy of its episode only if ADD-AUC rises by at least 0.03, neither success rate drops by more than 0.02, and MPPE rises by at most 0.5\,cm (analogous rules on Tier 1).
A new list is uploaded only if the mean ADD-AUC rises by at least 0.010 on the exported rollout and 0.005 on 4096 worlds, and if the predicted board score passes the next team by 0.002.
The prediction multiplies the smaller of the two gains by 0.21, the board to own-reference ratio measured on a Tier 2 upload (+0.0039 against +0.0187).

\paragraph{Tier 2 denoiser.}
A hand-designed filter tuned on Tier 2 data only holds static segments on a support surface, bridges motion segments with Brownian-bridge interpolation, fixes flips and height, removes penetration by the smallest rotation and smooths, before retargeting.

\paragraph{Tier 3 reconstruction.}
The CARI4D pipeline segments the video with Grounding DINO and SAM\,2~\cite{liu2024groundingdino,ravi2024sam2}, estimates depth with MoGe~\cite{wang2025moge}, tracks the object with FoundationPose~\cite{wen2024foundationpose} and fits the SOMA body, with a focal length of 866\,px from a ceiling-grid calibration (default 830\,px).
Static objects slide along their viewing ray onto the support, a dynamics pass removes implausible object motion, and object poses are picked among FoundationPose candidates for temporal consistency.
As the organizers allowed on 21 October, one rigid transform per episode aligns frame 0 with the Tier 1 frame-0 pose; no later Tier 1 frame is used.

\section{Experiments and Results}
\label{sec:experiments}

We trained about 370 policies in the official simulator on one RTX\,4090, one L40S and one L40.
Table~\ref{tab:board} gives the board scores. On Tier 3, frame-0 alignment, reconstruction fixes and per-episode swaps raised the board ADD-AUC from 0.106 to 0.282 and SP-SR from 0.267 to 0.467 between 18 and 29 October; on Tier 1, per-episode swaps lowered MPPE from 24.5\,cm (13 October) to 8.47\,cm.

\begin{table}[t]
  \centering
  \small
  \begin{tabular}{lcccc}
    \toprule
    Tier & ADD-AUC & SP-SR & MP-SR & MPPE (cm) \\
    \midrule
    1 & 0.635 & 0.867 & 0.700 & 8.47 \\
    2 & 0.518 & 0.733 & 0.533 & 15.32 \\
    3 & 0.282 & \textbf{0.467}$^{\ast}$ & 0.100 & 36.60 \\
    \bottomrule
  \end{tabular}
  \caption{Public board scores of maniphys.ai on 3 November 2026. $^{\ast}$First place on the board.}
  \label{tab:board}
\end{table}

\begin{table}[t]
  \centering
  \small
  \begin{tabular}{llccc}
    \toprule
    Tier & Ep. & Variant A & +\,\rtho & Swap \\
    \midrule
    2 & 0  & 0.400 & 0.774 & yes \\
    2 & 10 & 0.397 & 0.594 & yes \\
    2 & 18 & 0.521 & 0.855 & yes \\
    2 & 21 & 0.516 & 0.503 & no \\
    2 & 26 & 0.465 & 0.715 & yes \\
    3 & 0  & 0.528 & 0.795 & yes \\
    3 & 12 & 0.507 & 0.672 & no (MP-SR) \\
    3 & 18 & 0.626 & 0.826 & yes \\
    3 & 24 & 0.731 & 0.835 & yes \\
    3 & 26 & 0.579 & 0.812 & yes \\
    \bottomrule
  \end{tabular}
  \caption{4096-world ADD-AUC on the tier's own reference without and with \rtho (same seed, buffer and reference).}
  \label{tab:r3o}
\end{table}

\paragraph{Orientation term.}
This ablation isolates the term: each pair in Table~\ref{tab:r3o} differs only in Eq.~\ref{eq:ori}.
Over the 12 tracked Tier 2 and Tier 3 episodes we tested, ADD-AUC rose on 11, by 0.08 to 0.37, and 10 met all swap conditions.
On Tier 1 ep17 the base recipe had an object position error of 8.9\,cm but an orientation error of 42.0$^\circ$; with the term the orientation error fell to 22.7$^\circ$ and SP-SR rose from 0 to 0.932.
The term does not help episodes that lose the object, and it can destabilize training: on Tier 2 ep21 the 4\,M run lost MP-SR (0.965 $\rightarrow$ 0.316), and the 10\,M run of Tier 3 ep26 collapsed to zero success, so that episode keeps its 4\,M run.

\paragraph{Per-episode fixes on Tier 1.}
These single-change tests ask whether the lost Tier 1 episodes share one cause.
A larger grasp-closure weight (5 $\rightarrow$ 20) lowered the 4096-world MPPE of ep9 from 40.5 to 17.0\,cm, and an earlier reset when the palm leaves the reference (0.15 $\rightarrow$ 0.08\,m) lowered ep25 from 114.1 to 68.7\,cm; neither helped other lost episodes, so the causes look episode-specific.

\section{Limitations}
\label{sec:limits}

Policies are per episode and all results are in simulation.
Each arm uses one seed; on Tier 1, paired seeds of one recipe differed by up to 0.35 in MP-SR.
Episodes that lose the object soon after first contact stay unsolved (7 on Tier 1, 7 on Tier 2, at least 9 on Tier 3).
Own-reference gains overstate board gains, most on Tier 3, where 12 multi-phase successes against our reconstruction became 3 on the board.

\balance

\section{Conclusion}
\label{sec:conclusion}

A dense object-orientation term, accepted per episode only under fixed rules, raised ADD-AUC on 11 of 12 tracked Tier 2 and Tier 3 episodes at the cost of two environment variables and about two GPU-hours per policy, and new episodes need no retraining.

%%%%%%%%%%%%%%%%%%%%%%%%%%%%%%%%%%%%%%%%%%%%%%%%%%%%%%%%%%%%%%%%%%%%%%%%

%%% References, acknowledgments and appendices do not count toward
%%% the two-page limit.
\bibliographystyle{plainnat}
\bibliography{references}

%%%%%%%%%%%%%%%%%%%%%%%%%%%%%%%%%%%%%%%%%%%%%%%%%%%%%%%%%%%%%%%%%%%%%%%%

%%% Optional appendix: starts on a new page and does not count toward
%%% the two-page limit. Delete this part if you do not need it.
\end{document}
